\section{Corrosion Performance in Saline and Seawater Environments}
\label{sec:corrosion}

\subsection{Test environments and methods}

Almost all published AM-HEA corrosion data come from laboratory saline media---most commonly 3.5~wt\% NaCl or artificial seawater prepared to ASTM D1141 \cite{astm_d1141}---using potentiodynamic polarization and electrochemical impedance spectroscopy under the conventions of ASTM G3/G5/G61 \cite{astm_g3,astm_g5,astm_g61}. These tests rank alloys and passive-film quality well, but they are not seawater service: they omit biofilms, tidal wet--dry cycling, crevice geometries, flow, and multi-month timescales. Data from natural marine exposure exist for conventional AM alloys (e.g., SLM AlSi10Mg immersed in the Adriatic for 3--6 months \cite{am_alsi10mg_2026}), but equivalent long-term exposures of AM HEAs are, to our knowledge, absent from the literature. This is the single most important evidence gap in the field.

\subsection{AM versus cast: the homogeneity advantage}

The consistent finding across recent studies is that AM improves the chloride corrosion performance of single-phase and near-single-phase HEAs relative to cast counterparts. Mechanisms are twofold: (i) homogeneous elemental distribution removes micro-galvanic couples between segregated phases, and (ii) lower micro-porosity and finer structure reduce pit initiation sites. Reported results include LPBF CoCrFeMnNi with higher pitting resistance and polarization resistance than cast material \cite{han2025}, and SLM CoCrFeMnNi whose corrosion current density drops by roughly 58\% relative to the as-cast alloy, attributed to a more uniform composition and a more stable passive film \cite{han2025,cagirici2024}. In AlCoCrFeNi$_{2.1}$ EHEA, the as-built SLM condition outperforms both annealed and cast specimens in 3.5~wt\% NaCl for exactly these reasons \cite{slm_ehea_2025}. The caveat, again, is that these advantages are hostage to defect control: the same reviews flag that secondary phases, intermetallics, and segregation introduced by poor parameters or heat treatment can create localized corrosion sites \cite{cagirici2024}. The picture is not uniform: where SLM builds carried higher densities of pores, cracks, and dislocations, the \emph{cast} counterpart outperformed printed CoCrFeMnNi in 3.5~wt\% NaCl \cite{frontiers2020slmcast}---build quality, not chemistry alone, sets the ranking (Section~\ref{sec:a2}).

\subsection{AM HEAs against the marine benchmark: 316L}

The most direct marine-relevant comparison is LPBF CoCrFeMnNi against LPBF 316L. In chloride solution, as-built and HIPed CoCrFeMnNi exhibit roughly 2.2 times better corrosion performance than 316L, with stable pseudo-passive behavior and higher pitting potentials. In acidic solution the ranking inverts completely: 316L outperforms the HEA by more than an order of magnitude \cite{lpbf_cantor_2024}. The lesson for marine specification is that HEA superiority in seawater is real but \emph{environment-specific}; acidic, sulfide-rich, or creviced microenvironments (under deposits, inside crevices, beneath biofilms) cannot be assumed benign on the strength of neutral NaCl tests alone \cite{qiu2017}.

\subsection{Composition strategies for seawater passivity}

Several alloying strategies are emerging specifically for chloride service. Mo additions to EHEAs (AlCoCrFeNi$_{2.1}$Mo$_x$) simultaneously raise strength and pitting potential, with in-situ Mo-oxide/Cr-oxide composite films stabilizing passivity \cite{mo_ehea_2026}. Ti or Cu doping of LMD AlCoCrFeNi improves corrosion resistance over the undoped alloy, though secondary intermetallics must be suppressed \cite{cagirici2024}. Zr-containing HEAs have shown improved resistance over Zr-free variants \cite{cagirici2024}, and Nb--Mo--Ta--W refractory compositions designed for marine use beat 316L overall, albeit with deep pits nucleating at minor segregates in one study \cite{cagirici2024}. CoCrW-type alloys deposited by LPBF show passivation that evolves over weeks of immersion, with corrosion concentrated at melt-pool boundaries---a reminder that AM's own solidification pattern becomes part of the corrosion microstructure \cite{lpbf_coocrw2023}.

\subsection{Anisotropy, surface condition, and coupled degradation}

Three second-order effects are known from AM metals generally and are only beginning to be characterized in HEAs. \emph{Build orientation} changes the density of melt-pool boundaries and cell structures exposed at a surface, and corrosion rates in LPBF alloys vary with orientation \cite{han2025,cagirici2024}. \emph{Surface condition} dominates localized attack: for LPBF 316L, as-built rough surfaces are markedly more SCC-susceptible than machined ones, and cracks propagate along melt-pool boundaries \cite{lpbf_316l_scc2022}; machined AM surfaces can even exceed wrought behavior in pitting tests \cite{nass2025marine}. \emph{Coupled degradation}---corrosion fatigue, SCC, erosion-corrosion, and cavitation---is where marine components actually fail, yet AM-HEA data in these modes are essentially nonexistent; hydrogen permeation evidence for LPBF CoCrFeMnNi is one of the few positive indicators \cite{lpbf_cantor_2024}.

\subsection{Galvanic and biofouling interactions}

For DED-clad or WAAM-on-steel components, galvanic coupling between the HEA deposit and the substrate (or between an AM HEA part and adjacent conventional alloys) must be assessed; passive HEAs are typically noble relative to structural steel, so coating defects concentrate attack on the substrate side, and compatibility with cathodic protection regimes (e.g., offshore CP design practice) needs verification \cite{nass2025marine}. Biofouling interactions with HEA surfaces---which couple biological films to under-film chemistry and may accelerate localized attack---have not yet been studied in any systematic way. Table~\ref{tab:corrosion} summarizes the key corrosion evidence base.
